%% ----------------------------------------------------------------------------
% CVG SA/MA thesis template
%
% Created 03/08/2024 by Tobias Fischer
%% ----------------------------------------------------------------------------

\chapter{Related Work}
\label{chap:relatedwork}

\section{Visual(-Inertial) Odometry and Single-Session SLAM}

Polygraph is built on top of, rather than instead of, single-session visual and
visual-inertial SLAM. Systems such as ORB-SLAM3 \cite{campos2021orbslam3} and VINS-Fusion
\cite{qin2019vinsfusionlocal,qin2019vinsfusionglobal} track a single session's camera (optionally fused with IMU measurements) in real time, producing a locally accurate but drift-affected trajectory
together with intra-session loop closures. This thesis treats such a system's raw per-frame
pose stream as Polygraph's input odometry (\mysecref{sec:method-odometry}) and does not
modify it, but rather optimizes it. Polygraph's contribution begins where single-session tracking ends, at the problem of relating different sessions' independently tracked trajectories to one
another and alleviating drift. The incremental smoothing formulation backends of these systems use for pose-graph optimization (most directly iSAM2 \cite{kaess2012isam2}, built on the factor-graph
formulation underlying GTSAM) is also what Polygraph's own online graph solve
(\mysecref{sec:method-graph}) is built on.

\section{Visual Place Recognition for Loop Closure}

Deciding which pair of keyframes might correspond to the same physical place, before any geometric verification is attempted, is a visual place recognition (retrieval) problem. Each image is embedded in a global descriptor, and the nearest-neighbor search over previously seen descriptors proposes revisiting candidates. NetVLAD \cite{arandjelovic2016netvlad} established the learned aggregation approach on which this line of work is based. More recent descriptors, such as SALAD \cite{izquierdo2024salad}, BoQ \cite{alibey2024boq}, and MegaLoc \cite{berton2025megaloc} improve retrieval accuracy and robustness to change of viewpoint and appearance. Polygraph uses one such descriptor (\mysecref{sec:method-clustering}) to
cast votes between clusters of keyframes between sessions. Unlike single-session loop closure, where the retrieval index only ever needs to answer queries against one session's own past, Polygraph's retrieval index spans every other session concurrently and must scale as new sessions, not just new keyframes within a session, are added online.

\section{Learned Two-View Matching and Metric Pose Recovery}

A retrieval hit is only a hypothesis. Polygraph geometrically verifies each candidate keyframe pair before it is allowed to become a graph edge (\mysecref{sec:method-pose-estimation}). Sparse feature matching between the two images uses XFeat \cite{potje2024xfeat} for feature extraction and LighterGlue for matching, a lighter, faster retraining of LightGlue \cite{lindenberger2023lightglue} distributed as part of the XFeat project \cite{potje2024xfeat}. The overall recipe of sparse matching lifted to metric scale with a monocular depth prior matches 2GO's own default variant (``CL-Absolute'') \cite{lim2025twogo}. Classical and widely used alternatives to this matching step include SuperPoint \cite{detone2018superpoint} paired with SuperGlue \cite{sarlin2020superglue}, and, at the other extreme, dense two-view reconstruction methods such as MASt3R \cite{leroy2024mast3r} that recover a metric point map directly rather than only sparse
correspondences. Polygraph favors XFeat  \cite{potje2024xfeat} and LighterGlue \cite{lindenberger2023lightglue} for their speed, consistent with the lightweight, map-less design goal set out in \mychapterref{chap:intro}. However, as we will show later, other matching methods such as \cite{loftr} or \cite{leroy2024mast3r} could be used for our purpose. Because two-view matching alone only recovers relative pose up to an unknown scale, Polygraph lifts the matched keypoints to metric 3D with a monocular metric depth estimator, Metric3D \cite{yin2023metric3d}, and checks that the resulting metric estimate agrees with the scale-free triangulated one before accepting a match, avoiding the need for a second overlapping camera or a dense stereo/LiDAR sensor purely to resolve scale. At the pose-graph level, once candidate edges are admitted, Polygraph's final robust solve uses Graduated Non-Convexity with a truncated least-squares cost (GNC-TLS) \cite{yang2020gnc} to down-weight or reject edges that remain outliers even after the gating described in \mysecref{sec:method-gate}.

\section{Multi-Session and Lifelong Mapping}

Fusing more than one session's map is a long-standing problem beyond single-session SLAM. Systems that maintain dense or semi-dense multi-map representations, such as ORB-SLAM3's Atlas \cite{elvira2019orbslamatlasrobustaccuratemultimap} or the open, multi-session visual-inertial mapping framework maplab \cite{schneider2018maplab}, merge sessions by matching and aligning their maps directly, which gives strong geometric evidence for a match but ties the cost of fusion to the cost of maintaining dense geometry across every session. Other methods such as \cite{covins} or \cite{covins-g} build on top of general single-session SLAM and give them multi-session capabilities. However, these methods also utilize the map for pose estimation purposes. Polygraph instead follows the sparse, pose-graph-only line of work described next, trading some of that direct  geometric evidence for a representation whose cost scales with the number of keyframes and sessions, not with scene density.

\section{2GO}
\label{sec:relatedwork-2go}

The closest prior work to this thesis is 2GO \cite{lim2025twogo}, a single-session visual SLAM
pipeline that avoids maintaining any 3D map at all. Given one continuous capture, 2GO builds a
sparse, keyframed pose graph directly: long recordings are split into subsessions purely for
tractable processing, and loop closures between them are found and confirmed through learned
retrieval, sparse feature matching, and a monocular depth prior for metric scale, then stitched
back into a single trajectory. Polygraph's own per-pair verification
(\mysecref{sec:method-pose-estimation}) is built directly on this same recipe. 2GO has no
notion of fusing genuinely independent sessions, different agents, devices, or days, since it
was designed around a single trajectory; \mysecref{sec:experiments-2go-comparison} reports a
direct, hardware-matched comparison against its published KITTI numbers, including an open gap
in loop-closure recall that this thesis surfaces rather than resolves.

\section{Datasets for Multi-Session Evaluation}

KITTI \cite{geiger2012kitti} is a standard, widely used outdoor driving benchmark, even though it is intended for single-session experiments, some sessions do have real overlap. In this work we will exploit this fact and show that using this information boosts accuracy of the trajectory optimization. Oxford RobotCar \cite{maddern2017robotcar}, with the later RTK-based ground truth release \cite{maddern2020rtk}, provides repeated traversals of the same route recorded months apart under different weather, season, and lighting. The hard split draws a paired easy/hard session split from it, sharing all but one session between the two, with the hard split adding a night traversal, specifically to test single- and multi-session performance under real day/night appearance change, axis prior work on this dataset has approached from other angles, including radar-based SLAM \cite{radarslam2021}, direct sparse odometry with omnidirectional cameras \cite{matsuki2018omnidso}, and a systematic evaluation of visual odometry robustness in rain \cite{tan2023voautorain}, but never for multi-session pose-graph fusion. The data sets evaluated in this thesis are built from existing localization/mapping benchmarks rather than newly captured: HGE (both a HoloLens and an iPad capture) from LaMAR \cite{sarlin2022lamar}, a benchmark for AR localization and mapping across long, unconstrained indoor/outdoor trajectories, and HYDRO (a Boston Dynamics Spot capture) from CroCoDL \cite{blum2025crocodl}, a cross-device dataset combining HoloLens and legged-robot recordings; neither has previously been used to evaluate multi-session pose graph fusion. Both are introduced together with the rest of the experimental setup in \mychapterref{chap:experiments}. EuRoC \cite{burri2016euroc} is also part of the larger experimental campaign (\mysecref{sec:experiments-euroc}). Here we will compare the anchor quality between sessions and the general pose estimation quality against existing methods such as \cite{elvira2019orbslamatlasrobustaccuratemultimap, covins, covins-g}. In the end, similar to \cite{lim2025twogo}, we will perform single and multi session test on large scale outdoor VBR Rome \cite{vbr} dataset.